% !TEX program = xelatex
% Spatial block-bootstrap sensitivity: does the significant ozone plain-CRPS
% cell survive once between-site spatial correlation is respected?
% Table generated by make_block_table.R -> block_table_body.tex.
\documentclass[12pt]{article}

\usepackage{fontspec}
\usepackage{xeCJK}
\setCJKmainfont[AutoFakeSlant=.2, AutoFakeBold=3]{AR PL KaitiM Big5}

\usepackage[margin=1in]{geometry}
\usepackage{amsmath, amssymb}
\usepackage{booktabs}
\usepackage{array}
\usepackage{hyperref}
\hypersetup{colorlinks=true, linkcolor=blue!45!black, urlcolor=blue!45!black}

\title{\textbf{站間相依敏感度：唯一的顯著格撐不撐得住？}}
\author{site bootstrap vs 空間 block bootstrap（回應文獻檢視 \#1）}
\date{2026 年 7 月}

\begin{document}
\maketitle

\section*{為什麼要做這個}

前面所有 CI 都用 \textbf{site-clustered bootstrap}：對驗證\emph{站}有放回重抽，把每
站（連同其所有日）當獨立 cluster。這正確處理了\emph{站內}（跨日）相依，卻假設
\textbf{站間空間獨立}。但臭氧鄰近測站傾向同日一起極端（skew-t 過程刻意建的就是
這個），所以有效 cluster 數 $<n_{\mathrm{site}}$，site-bootstrap 的 CI \textbf{偏窄
（anticonservative）}。Cressie 對相依資料盲用 bootstrap 的警告、Wikle 等的
effective d.o.f.\ 論點，都指向這一點。

\textbf{檢驗方式}：把域切成連續地理區塊，改對\emph{區塊}重抽（鄰近站一起帶走）；
區塊愈粗，吸收愈長程的相依、CI 愈寬。對 null 而言偏窄 CI 是安全方向；真正有風險的
是唯一的顯著格——臭氧 plain-CRPS 的「AR(2) 較差」。

\bigskip
% Auto-generated by make_block_table.R -- do not edit by hand.
\begin{table}[htbp]
\centering
\caption{Spatial block-bootstrap sensitivity, ozone AR(2) vs.\ AR(1). $p$ is
the two-sided bootstrap $p$ for $\Delta=\mathrm{CRPS}_{AR(2)}-\mathrm{CRPS}_{AR(1)}$;
a star marks a $95\%$ CI excluding $0$. Columns move from resampling single
sites (assumes spatial independence) to resampling geographic blocks (keeps
spatially-correlated neighbours together; coarser blocks absorb longer-range
correlation). The two tail scores stay null throughout; the plain-CRPS
``AR(2) worse'' is significant only under the (anticonservative) site
resampling and dissolves as soon as blocks respect spatial dependence.}
\label{tab:block-sensitivity}
\small
\begin{tabular}{@{}ll c cccc@{}}
\toprule
Split & Score & $\Delta$ & site & block (12) & block ($\approx$20) & block ($\approx$32) \\
\midrule
\multicolumn{7}{@{}l}{\emph{200/200 split}}\\
& Brier $q_{.95}$ & $-0.000062$ & $0.753$ & $0.696$ & $0.771$ & $0.758$ \\
& twCRPS (tail) & $-0.00015$ & $0.981$ & $0.959$ & $0.993$ & $0.951$ \\
& plain CRPS & $+0.021$ & $0.026^{*}$ & $0.113$ & $0.082$ & $0.076$ \\
\addlinespace
\multicolumn{7}{@{}l}{\emph{300/100 split}}\\
& Brier $q_{.95}$ & $-0.000027$ & $0.808$ & $0.834$ & $0.825$ & $0.912$ \\
& twCRPS (tail) & $-0.00021$ & $0.868$ & $0.913$ & $0.899$ & $0.935$ \\
& plain CRPS & $+0.039$ & $0.027^{*}$ & $0.109$ & $0.091$ & $0.047^{*}$ \\
\bottomrule
\end{tabular}
\end{table}


\section*{結果：null 全穩健，唯一的顯著格瓦解}

\paragraph{兩個尾部分數（Brier、twCRPS）：處處 null，穩健。}從 site 到最粗的區塊，
$p$ 值幾乎不動（Brier $\approx0.75$--$0.91$、twCRPS $\approx0.87$--$0.99$），CI 只是
適度變寬仍含 $0$。這證實了文獻檢視的判斷：\textbf{對 null，偏窄的 site CI 是安全的}
——更保守的區塊 bootstrap 只會讓 null 更穩。

\paragraph{plain CRPS 的「AR(2) 較差」：撐不住。}它在 site 重抽下顯著
（$p=0.026$／$0.027$，兩個 split），但一旦區塊尊重空間相依，$p$ 就升到
$0.08$--$0.11$、CI 含 $0$（$200/200$ 全部 ns；$300/100$ 僅最細網格勉強 $0.047$）。
換言之，\textbf{那個「顯著、可複現」是忽略站間相依造成的假象}；點估計沒變，但它
不再與 $0$ 可分。

\section*{對先前結論的（第二次）修正}

folder 05/06 我把臭氧 plain-CRPS 的「AR(2) 較差」當作一個\emph{真實的} bulk 效應
（雖已由 twCRPS 判定離題）。現在要再退一步：\textbf{它連統計上都不穩健}。所以這條
線的完整、誠實版本是：

\begin{quote}
臭氧 plain-CRPS 上「AR(2) 較差」有\emph{兩重}問題：（一）它是 bulk 效應，不在尾部
（folder 06/07 的 twCRPS）；（二）它對站間空間相依不穩健，空間 block bootstrap 下
不再顯著（本文）。無論從\emph{目標對口}或\emph{統計穩健}任一角度，它都不承載結論。
\end{quote}

\paragraph{淨效果：主結論更強。}論文的核心結論全是 \textbf{null}（AR(2)/MRTS 不改變
尾部超標預測）。這些 null 現在通過了空間 block bootstrap 的加壓測試——CI 變寬仍含
$0$。而唯一曾經「顯著」的格被證明是假象。整條不確定性分析至此密不透風：\textbf{兩項
延伸在任何情境、任何尾部分數、且計入空間相依後，都不改變門檻超標的預測技巧。}

\medskip
\noindent\footnotesize 程式：\verb|ozone/US-all-auto/block_bootstrap_sensitivity.R|
（地理格網區塊、ratio 估計量、$B{=}2000$）。區塊數依網格與非空格而定
（粗 $\approx12$、中 $\approx18$--$20$、細 $\approx28$--$35$）。
\end{document}
